\section{有限七次被覆と分岐曲線の継続入口}
\label{sec:notes-septic-new-v30}

\paragraph{一般点で確定した幾何.}
選択した十次元余因子成分 $Y_0$ の非空開部分で
\[
 f:X=\operatorname{Bl}_{18}\mathbb P^2\longrightarrow\mathbb P^2,
 \qquad L=5H-\sum E_i
\]
は有限平坦な次数七の射であり、$L$ は豊富である。
一般の幾何学的monodromyは $S_7$ である。
分岐源 $R$ は滑らかな幾何学的整曲線で、
\[
 R\sim12H-2\sum E_i,\quad R^2=72,\quad K_XR=0,\quad g(R)=37.
\]
被約分岐曲線 $B$ は整な二十四次曲線で、$\nu:R\to B$ は正規化となる。
従って $p_a(B)=253$、$\sum\delta_p(B)=216$ である。
微分 $d\nu$ の零因子の長さは七十二、双対曲線 $B^*$ の次数は四十八である。
公表稿VIIではこれらを一般の同じ成分上で証明する\cite{PaperVII}。

\paragraph{有限性の証明に再利用できる判定.}
$h^0(\eta(H))=0$、$h^0(\eta(2H))=3$ と後者の基点自由性から
$h^0(\eta(H+p))=0$ を得る。従って十八点のうち任意の十七点を通る四次曲線はない。
$L$ が縮約する既約曲線を $dH-\sum m_iE_i$ とすると
\[
 \sum m_i=5d,\qquad \sum m_i^2\le d^2+2d+2.
\]
Cauchyから $d\le6$、整数の最小平方和から $d=5,6$ を除き、
上の十七点条件から $d=4$ も除く。
残るのは五点直線、十点二次曲線、十五点三次曲線、
または一点二重と十三点単純の三次曲線である。
同じ成分上の六変数特殊族では、縮約曲線は三本の $F_i=0$ だけである。
十八点を標識した有限エタール被覆上で三つの効果的条件が真閉集合となることを
一次変形で示せば、一般点の有限性が従う。
標識被覆の全成分を一つと見なさず、証人を通る既約成分を選んでから開集合を降ろす。

\paragraph{分岐源の整性を証明する特殊化.}
単独foldの証人 $A_*=(3,-1,4,-2,5/4,3/4)$ では、残余critical六次式は二次式と四次式の積に分解する。
この点を滑らかな残余曲線の証人として再利用してはならない。
六変数族の点 $(a,c,h,j,n,q)=(1,2,3,4,5,6)$ では
\[
 R_*=R_0+T_0+T_1+T_2.
\]
$R_0$ は滑らかな種数十の六次曲線のstrict transform、
$T_i$ は互いに交わらない収縮二次曲線のstrict transformである。
全成分の係数は一、$LR_0=24$、$LT_i=0$ となる。
一般点では $L$ が豊富なので、一般分岐源の各成分の特殊化には
唯一の正次数成分 $R_0$ が必要になる。その係数が一だから一般分岐源は幾何学的に整である。
単独foldの持続と合わせ、分岐源から被約分岐曲線への次数が一になる。
この特殊族自体の分岐曲線の正規化種数十を、一般点の種数と同一視しない。

\paragraph{一次で見えない節点を平滑化する二次補正.}
$R_*$ は三十の通常節点を持ち、そのうち二十四点はblow-up基点の外にあり、
六点は例外曲線上にある。前者では一次平滑化が検出できる。
後者では動く基点を二次まで追い、blow-up座標で実際の分岐式を取る必要がある。
局所分岐式の二次部分を
\[
 r(z,\varepsilon)=\tfrac12z^tHz+\varepsilon c^tz+\varepsilon^2b+
 \text{三次以上},\qquad \det H\ne0
\]
と書くと、真の平滑化係数は $b$ 単独ではなく
\[
 b-\tfrac12c^tH^{-1}c
\]
である。これは臨界点 $z=-\varepsilon H^{-1}c+\cdots$ を代入した値である。
六点でこの値は $\mathbb F_{163}$ 上
$118,118,14,14,31,31$ となり、全て非零である。
動く基点に対し $h=g/u$ として $\det[h;h_u;h_w]$ を計算する方法と、
斉次Jacobianを $5u^2$ で割る方法が独立に一致する。
係数の $\varepsilon^2$ 項を二階微分そのものと取り違えない。

標数零へ戻すには、孤立した有限体の飽和計算だけでは不十分である。
三十節点の交差族が固有かつ有限エタールとなる局所環を取り、
例外上六点を閉開部分として区別する。
単元Hessianと非零一次・二次係数を持ち上げ、相対特異点スキームの固有性を使えば、
有限個の真の判別条件を避けた一般点で $R$ が滑らかとなる。

\paragraph{七十二と四十八を特異点型なしで使う.}
$R$ が滑らかなので $df$ は $R$ 上で階数一であり、
$\operatorname{coker}(df)=i_*\mathcal M$ で $\mathcal M$ は線束となる。
曲面Riemann--Rochから
\[
 \chi(T_X)=-28,\quad\chi(f^*T_{\mathbb P^2})=8,
 \quad\deg\mathcal M=72.
\]
$\mathcal I=\operatorname{im}(df|_R)$ の次数は $3LR-72=0$ なので、
$T_R\to\mathcal I$ の零因子 $Z_\nu$ は次数七十二である。
Gauss写像の線束は
$K_R\otimes\nu^*\mathcal O(2)(-Z_\nu)$ で、次数は四十八となる。
標数零の双対性によりGauss写像は双有理である。

尖点七十二個・節点百四十四個という個数は、分岐曲線が通常の節点と尖点だけを持つ場合の帰結である。
次の有効な検証対象は、$Z_\nu$ の被約性、分岐源の二点を同じ像へ送る対応の横断性、
三重像と非通常接触の排除である。これらを七次元の制約付き族の中で証明する必要があり、
無制約な一般の曲面射の安定性をそのまま引用して済ませない。
